\subsection*{MT-M25. Exact line-completion kernel and the strict jet cutoff}
Fix the Hecke function, $\mathcal E$, $\Xi=\Xi_\eta$,
$\mathcal I=\mathcal I_\eta$ and $Q=\mathcal E/\mathcal I$
of MT-M18. For $s\ge0$ and $a\in\mathbf R$, put
\[
 H_s=L^2\bigl(\mathbf R,(1+x^2)^{s/2}\,dx\bigr),
 \qquad T_af(x)=e^{ax}f(x),
 \qquad J_{a,s}=\overline{T_a\mathcal I}^{\,H_s},
 \quad Z_{a,s}=H_s/J_{a,s}.
 \tag{M25.1}
\]
These are exactly the fixed-character versions of Masullo's
line norm (9.1.2.3), original PDF p.~399, his closure
(10.1.1.3), p.~434 with $I=\{0\}$, and his quotient
(10.2.1.2), p.~439, by MT-M24. The fixed unitary norm
modulation has absolute value one and is harmless for this norm.
Let $m_\rho$ be the multiplicity of a zero of $\Xi$, and set
\[
       q_s=\max\{0,\lceil(s-1)/2\rceil\}.
 \tag{M25.2}
\]
Then
\[
 J_{a,s}=
 \bigcap_{\substack{\rho=a+i\gamma:\Xi(\rho)=0\\
                    0\le k<\min(m_\rho,q_s)}}
 \ker\left[u\longmapsto\int_{\mathbf R}
                  u(x)x^k e^{i\gamma x}\,dx\right].
 \tag{M25.3}
\]
Every functional displayed here is continuous on $H_s$.
The natural map
\[
 q_{a,s}:Q\longrightarrow Z_{a,s},\qquad[f]\longmapsto[T_af]
 \tag{M25.4}
\]
is continuous with dense image, and its kernel consists exactly
of the classes whose jets $F^{(k)}(\rho)$, $F=\mathcal Mf$,
vanish at all the indices in (M25.3). Thus the line completion
retains only the zeros on $\Re\rho=a$ and only their derivatives
with the \emph{strict} inequality
\[
                         2k<s-1.
 \tag{M25.5}
\]
It is faithful on the full smooth quotient if and only if all
zeros lie on that line and $m_\rho\le q_s$ for every zero.
No GRH or simplicity is assumed in these assertions.

\begin{proof}
The continuous bilinear dual of $H_s$ can be written as
$H_{-s}$ under $\ell_g(u)=\int u(x)g(x)\,dx$.
This is just the Riesz representation theorem after multiplying
by $(1+x^2)^{s/4}$; complex conjugation can be absorbed into
the representing $g$.
For a polynomial $P$ of degree $k$ with nonzero leading
coefficient,
\[
 P(x)e^{i\gamma x}\in H_{-s}
 \quad\Longleftrightarrow\quad
 \int_{|x|\ge1}|x|^{2k-s}\,dx<\infty
 \quad\Longleftrightarrow\quad 2k<s-1.
 \tag{M25.6}
\]
This proves continuity of the displayed functionals and the
necessity of their strict numerical threshold. It also shows
that the number of admissible nonnegative integers $k$ is
exactly $q_s$, including $0\le s\le1$ when there are none.
Each such functional vanishes on $T_a\mathcal I$ by MT-M18's
full-jet ideal equality, hence on its Hilbert closure.

We prove there are no other annihilators. The elementary
inequality
\[
 1+(x+y)^2\le2(1+x^2)(1+y^2)
\]
implies, for either sign of $s$,
\[
 \|u(\,\cdot-y)\|_{H_{\pm s}}
       \le 2^{s/4}(1+y^2)^{s/4}\|u\|_{H_{\pm s}}.
 \tag{M25.7}
\]
For the negative sign use the same inequality with $x$ and
$x-y$ interchanged before taking reciprocals. Minkowski's
integral inequality therefore gives
\[
 \|k*u\|_{H_{\pm s}}
 \le2^{s/4}\|u\|_{H_{\pm s}}
       \int|k(y)|(1+y^2)^{s/4}\,dy.
 \tag{M25.8}
\]
In particular convolution by a Schwartz function is bounded
on both spaces. Smooth compactly supported functions are
dense in both weighted $L^2$ spaces: truncate first, and then
approximate on the remaining compact interval, where the
weight and its inverse are bounded. These facts also prove
strong continuity of translation, first on that dense subspace
and then by (M25.7).

Let $g\in H_{-s}$ annihilate $T_a\mathcal I$.
It defines a tempered distribution, since Cauchy--Schwarz
bounds $|\int ug|$ by the continuous Schwartz seminorm
$\|u\|_{H_s}\|g\|_{H_{-s}}$.
Use the Fourier convention
$\mathcal F_+u(t)=\int u(x)e^{itx}\,dx$ and its inverse
$(2\pi)^{-1}\int\phi(t)e^{-itx}\,dt$.
Define the tempered distribution $D$ by
$D(\phi)=\ell_g(\mathcal F_+^{-1}\phi)$.
Since $\psi_a=T_a\psi_\eta\in\mathcal E$ and
$\mathcal F_+\psi_a(t)=\Xi(a+it)$, we obtain
\[
                         \Xi(a+it)D=0.
 \tag{M25.9}
\]
Here is why the identity holds on every Schwartz test,
not just an entire Fourier transform. For $u\in\mathcal E$,
$\psi_a*u\in T_a\mathcal I$ because
$T_a(f_1*f_2)=T_af_1*T_af_2$ and $T_a\mathcal E=\mathcal E$.
Thus $\ell_g(\psi_a*u)=0$. The space $\mathcal E$ contains
$C_c^\infty$ and is dense in the usual Schwartz space;
convolution by $\psi_a$ is continuous there, and
$\ell_g$ is continuous there. This extends the equality to
all Schwartz $u$, which is precisely (M25.9).

The zeros of $t\mapsto\Xi(a+it)$ on the real line form a
locally finite set. Away from them, a compactly supported
smooth test can be divided by this nonvanishing multiplier,
so (M25.9) shows $\operatorname{supp}D$ is contained in that
set. Near a zero $\gamma$ of order $m_\rho$, multiplication
by a smooth invertible function reduces (M25.9) to
$(t-\gamma)^{m_\rho}D=0$. A distribution supported there
with this property depends only on the first $m_\rho$ jets
of a test: subtract its Taylor polynomial of degree
$m_\rho-1$ times a cutoff equal to one near $\gamma$;
the remainder is divisible by $(t-\gamma)^{m_\rho}$ there.
It is therefore a linear combination of
$\delta_\gamma^{(k)}$, $0\le k<m_\rho$.

Choose a smooth frequency cutoff supported near this single
zero and equal to one near it. Multiplying $D$ by that
cutoff corresponds to convolution of $g$ with a Schwartz
kernel, with an inessential reflection dictated by the
chosen Fourier convention. By (M25.8) the resulting
function still belongs to $H_{-s}$. Its inverse transform
is a polynomial times $e^{i\gamma x}$. Formula (M25.6)
therefore excludes every nonzero coefficient of degree
$k$ with $2k\ge s-1$. Locally only the functionals listed
in (M25.3) occur.

Finally these local pieces exhaust the annihilator in its
Hilbert norm. Choose $\vartheta\in C_c^\infty(\mathbf R)$
equal to one near zero and put
$\vartheta_R(t)=\vartheta(t/R)$, $R\ge1$.
The associated convolution kernels have the form
$Rk(Rx)$, possibly reflected, where $k$ is Schwartz and
$\int k=1$. Their weighted $L^1$ norms in (M25.8) are
uniformly bounded, because
\[
 \int R|k(Ry)|(1+y^2)^{s/4}\,dy
       =\int|k(v)|(1+v^2/R^2)^{s/4}\,dv
       \le\int|k(v)|(1+v^2)^{s/4}\,dv.
\]
Translation continuity and dominated convergence in the
same integral show $Rk(R\cdot)*g\to g$ in $H_{-s}$.
The distribution $\vartheta_RD$ is supported on finitely
many zeros, so each such convolution is a finite linear
combination of the permitted polynomial waves. Hence
the full annihilator of $T_a\mathcal I$ is the Hilbert
closure of the span of precisely the functionals in
(M25.3). Taking its annihilator proves (M25.3).

The inclusion $T_a:\mathcal E\to H_s$ is continuous,
since one exponential seminorm with weight greater than
$|a|$ absorbs the polynomial factor in (M25.1). Its range
is dense because $T_a\mathcal E=\mathcal E$ contains
$C_c^\infty$. It vanishes on $\mathcal I$ after passing
to the quotient, and so induces (M25.4), continuously
with dense image. The functional indexed by $(\rho,k)$
evaluates $T_af$ to
\[
       \int e^{ax}f(x)x^ke^{i\gamma x}\,dx
                    =(\mathcal Mf)^{(k)}(\rho).
\]
This gives its exact kernel. If all zero jets are retained,
MT-M18 makes that kernel zero. If any zero jet is omitted,
the finite-block interpolation of MT-M18 realizes a class
with just that jet nonzero and all other zero jets zero.
It is a nonzero class in the kernel of (M25.4). This proves
the stated necessary and sufficient condition for faithfulness.
\end{proof}

\paragraph{A literal boundary correction with an arithmetic witness.}
The original PDF (10.2.1.6), p.~440, displays
\[
 N_s(\rho)=\min\{N(\rho),[(s-1)/2+1]\}.
\]
At $s=3$ this allows the first derivative at any double zero.
But its proposed representing function is
$x e^{i\gamma x}$, whose squared $H_{-3}$ norm contains
\[
              \int_{|x|\ge1}\frac{x^2}{(1+x^2)^{3/2}}\,dx
                         =+\infty.
 \tag{M25.10}
\]
Thus the endpoint inclusion is false as written; the
correct retained multiplicity is
$\min\{m_\rho,q_s\}$. Both the source norm (9.1.2.3),
p.~399, and the displayed cutoff (10.2.1.6), p.~440, were
rendered and visually read from the original PDF.
This is not a hypothetical objection requiring an unknown
double zero of the Riemann zeta function. By NT-L17,
the actual degree-six field
$K=\mathbf Q(\sqrt[3]{2},\sqrt{-3})$ has nontrivial zeros
of $\zeta_K$ of order at least two. Choose such a zero
$\rho$, put $a=\Re\rho$, and use its trivial Hecke
character. Then (M25.10) is a concrete counterexample
to the stated membership at its permitted $a$ and $s=3$.
Under $M$, one may take $a=1/2$.

More generally, for every fixed finite $s$, choose an odd
prime $r$ with $r-1>q_s$. The Kummer field of NT-L17
has zeros of multiplicity at least $r-1$. Its line quotient
at $a=1/2$ cannot faithfully retain the full smooth
quotient: under $M$ it loses the higher jets, and without
$M$ either that loss occurs or an off-line zero is lost.
The strict cutoff is a property of this particular
Hilbert completion, not a failure of the unconditional
adapted-block synthesis in MT-M18.

This gives the precise analytic map at the inherited
cote~1 / SA23 geo-H-lidskii and geo-H-genus interface:
the Hilbert norm can discard part of the finite-length
spectral divisor. It does not prove or refute the
conditional endpoint $M$, and it invokes none of the
draft's unverified general basis or nuclearity arguments.
